\documentclass[letterpaper,11pt]{article}

\usepackage{latexsym}
\usepackage[empty]{fullpage}
\usepackage{titlesec}
\usepackage{marvosym}
\usepackage[usenames,dvipsnames]{color}
\usepackage{verbatim}
\usepackage{enumitem}
\usepackage[pdftex]{hyperref}
\usepackage{fancyhdr}
\usepackage{graphicx}
\usepackage{fontawesome}

\pagestyle{fancy}
\fancyhf{} % clear all header and footer fields
\fancyfoot{}
\renewcommand{\headrulewidth}{0pt}
\renewcommand{\footrulewidth}{0pt}

% Adjust margins
\addtolength{\oddsidemargin}{-0.375in}
\addtolength{\evensidemargin}{-0.375in}
\addtolength{\textwidth}{1in}
\addtolength{\topmargin}{-.5in}
\addtolength{\textheight}{1.0in}

\urlstyle{same}

\raggedbottom
\raggedright
\setlength{\tabcolsep}{0in}

% Sections formatting
\titleformat{\section}{
  \vspace{-4pt}\scshape\raggedright\large
}{}{0em}{}[\color{black}\titlerule \vspace{-5pt}]

%-------------------------
% Custom commands
\newcommand{\resumeItem}[2]{
  \item\small{
    \textbf{#1}{: #2 \vspace{-2pt}}
  }
}

\newcommand{\resumeSubheading}[4]{
  \vspace{-1pt}\item
    \begin{tabular*}{0.97\textwidth}{l@{\extracolsep{\fill}}r}
      \textbf{#1} & #2 \\
      \textit{\small#3} & \textit{\small #4} \\
    \end{tabular*}\vspace{-5pt}
}

\newcommand{\resumeSubItem}[2]{\resumeItem{#1}{#2}\vspace{-4pt}}

\renewcommand{\labelitemii}{$\circ$}

\newcommand{\resumeSubHeadingListStart}{\begin{itemize}[leftmargin=*]}
\newcommand{\resumeSubHeadingListEnd}{\end{itemize}}
\newcommand{\resumeItemListStart}{\begin{itemize}}
\newcommand{\resumeItemListEnd}{\end{itemize}\vspace{-5pt}}

%-------------------------------------------
%%%%%%  CV STARTS HERE  %%%%%%%%%%%%%%%%%%%%%%%%%%%%


\begin{document}

%----------HEADING-----------------
\begin{tabular*}{\textwidth}{l@{\extracolsep{\fill}}r}
  \textbf{\Large Surya Teja Eada} & Email : \href{mailto:surya.eada@uconn.edu}{suryateja.eada@oregonstate.edu}\\
Department of Statistics   & Mobile : +1 (860) 771-8215 \\
Oregon State University & \href{https://scholar.google.com/citations?user=gxKgkdgAAAAJ\&hl=en\&oi=ao}{Google Scholar}
\end{tabular*}

%-----------EXPERIENCE-----------------
\section{Experience}
  \resumeSubHeadingListStart

   \resumeSubheading
     {Oregon State University}{Corvallis, Oregon}
     {Assistant Professor of Teaching}{Sep. 2026 -- Present}	
     
    \resumeSubheading
     {Truist Financial}{Charlotte, North Carolina}
     {Senior Model Development Analyst}{Apr. 2025 -- Jul. 2026}	

   \resumeSubheading
     {Oregon State University}{Corvallis, Oregon}
     {Assistant Professor of Teaching}{Sep. 2024 -- Mar. 2025}	

    \resumeSubheading
    {University of Connecticut}{Storrs, CT}
    {Graduate Teaching Assistant, Dept. of Statistics and Dept. of Mathematics}{Aug. 2017 - Aug. 2024}

    \resumeSubheading
      {CRISIL}{Pune, India}
      {Senior Quantitative Analyst}{Jul. 2014 - May 2017}
  \resumeSubHeadingListEnd



%-----------EDUCATION-----------------
\section{Education}
  \resumeSubHeadingListStart
    \resumeSubheading
      {University of Connecticut}{Storrs, CT}
      {Ph.D. in Statistics; GPA: 4.09/4}{Aug. 2019 -- Aug. 2024} \\ 
      \begin{itemize}
      \item Dissertation: \textit{L\'evy Processes
      driven by Telegraph Signal Process}
      \item Major Co-advisors: Vladimir Pozdnyakov and Jun Yan
      \end{itemize}
    \resumeSubheading
      {University of Connecticut}{Storrs, CT}
      {Master of Science in Applied Financial Mathematics;  GPA: 4.05/4}{Aug. 2017 -- May 2019}
    \resumeSubheading
      {Indian Statistical Institute}{Kolkata, India}
      {Master of Statistics }{Jul. 2012 -- May 2014}
    \resumeSubheading
      {Indian Statistical Institute}{Kolkata, India}
      {Bachelor of Statistics }{Jul. 2009 -- May 2012}
  \resumeSubHeadingListEnd

\section{Research Areas}
\begin{itemize}
\item Stochastic Processes, Hidden Markov Models, Parametric Modeling, Adversarial Machine Learning, AI
\end{itemize}

\section{Research Work}
\begin{itemize}

\item L\'evy Processes Governed by Telegraph Signal: Parametric Modeling and Numerical Estimation using L\'evy processes such as Brownian process, Meixner process to describe movement while underlying state is governed by a telegraph signal. Code implemented in R-cpp and R. Ongoing work with a paper published in \textit{Statistical Inference for Stochastic Processes}.

\item Adversarial Machine Learning: Parametric Modeling of Success/Failure of Adversarial Machine Learning examples against defense ensembles or real world systems such as elections to then do Estimation of Parameters that represent relevant attributes. Ongoing work.

\end{itemize}

%-----------Publications-----------------

\section*{Publications (\href{https://scholar.google.com/citations?user=gxKgkdgAAAAJ\&hl=en\&oi=ao}{Google Scholar})}
\textbf{Eada, S. T.}, Pozdnyakov, V., \& Yan, J. (2025). Discretely observed Brownian motion governed by telegraph signal process: Estimation and application to finance. \textit{Statistical Inference for Stochastic Processes}, 28(1), 1.

\section*{E-Prints (not peer reviewed)}
Kamal, K. M. A., \textbf{Eada, S.}, Verma, A., Acharya, S., Yemin, A., Fuller, B., \& Mahmood, K. (2026). Analyzing Physical Adversarial Example Threats to Machine Learning in Election Systems. arXiv preprint arXiv:2603.00481, \href{https://arxiv.org/pdf/2603.00481}{Eprint}

\section*{Presentations}

\begin{itemize}
    \item Delivered an invited talk on ``Brownian Processes driven by
    Telegrapher's process'' at the New England Statistical Symposium (NESS) 2023.
    \item Delivered a workshop ``Introduction to Python'' at
    University of Connecticut Sports Analytics Symposium (UCSAS) 2021.
\end{itemize}

\section*{Awards}

\begin{itemize}
    \item H. Fairfield Smith and Dolores S. Smith Award for
    Excellence in Applied Statistics, University of Connecticut (2020)  
    \item Certificate of Excellence for the Best Performance in
    Probability, University of Connecticut (2021)
\end{itemize}  

\section{Work Experience}
  \resumeSubHeadingListStart
     \resumeSubheading
    {Oregon State University}{Corvallis, OR}
    {Assistant Professor of Teaching, Dept. of Statistics}{Sep. 2024 - Mar. 2025, Sep. 2026 - Present}
      \begin{itemize}
          \item Taught Sampling Methods (ST 431/531),  Methods of Data Analysis (ST 412/512) and 
          Time Series Analytics (ST 566)
          \item Taught across both In-person and Online modalities
          \item Taught slash courses that cater to both undergraduate and graduate students
          \item Served on Graduate Admission Committee
      \end{itemize}
      
     \resumeSubheading
    {Truist Financial}{Charlotte, NC}
    {Senior Model Development Analyst}{Apr. 2025 - Jul. 2026}
      \begin{itemize}
          \item Worked in a model development role running exploratory data analysis, building credit loss
          forecasting models for the Comprehensive
        Capital Analysis and Review (CCAR) exercise in particular for the prime and sub-prime auto portfolios.  
      \end{itemize}

    \resumeSubheading
    {University of Connecticut}{Storrs, CT}
    {Teaching Assistant, Dept. of Statistics and Dept. of Mathematics}{Aug. 2017 - Aug.2024}
      \begin{itemize}
          \item Accumulated a wealth of teaching experience over 14 semesters.
          \item Held diverse roles including leading discussions,
           instructing a complete online course, and serving as the primary instructor for courses.
          \item Taught across different levels (undergraduate/graduate) within both statistics and mathematics departments.
          \item Served on Graduate Student Committee and was President of TARANG (South Asian Graduate Student Committee)
      \end{itemize}

    \resumeSubheading
      {CRISIL}{Pune, India}
      {Senior Quantitative Analyst}{July 2014 - May 2017}
      \begin{itemize}
        \item Credit Risk Modeling: As an outsourcing consultant for a US commercial bank,
        I specialized in model validation working for Comprehensive
        Capital Analysis and Review (CCAR) stress testing models. Here, I
        validated logistic regression models utilized for Probability
        of Default and fractional outcome regression models for Loss
        Given Default for HELOC and Auto loan portfolios. I also helped in responding to questions on validation
         during federal audit and excelled in a mentoring role.
                 
        \item Operation Risk Models: Detected inappropriate use of Poisson distribution,
            preventing underestimation of operational loss values for
            specific events.
        \item House pricing models: Offered feedback in the inaccurate use of type of
        testing (Wilcoxon test instead of Wilcoxon signed rank
            test) that cut client spending, resulting in cost-saving.
      
    \end{itemize}


    \vspace{-1pt}\item
      \textbf{Internships}
    \begin{itemize}
        \item Institute for Mathematical and Statistical
        Innovation (IMSI), Storrs, CT, Summer 2021: Worked on Leaf Area
        and Yield Model where we Utilized a partitioning logistic module to
        study sorghum, evaluating yield by leaf weight per unit area.
        Determined actual yields and areas for optimal model
        parameters, comparing model predictions with real-world data.
        Explored various loss definitions, ultimately adopting a
        normed Mahalanobis loss to optimize parameters in the module.
        
        \item Traveler's, Hartford, CT, Spring 2019: Collaborated on a
        semester-long experiential project to assess the impact of
        modern technology and direct-to-consumer channels on Traveler's
        primary distribution method, the "Independent Agency channel"
        and provide recommendations. 
        \item Deloitte U.S. India Ltd., Hyderabad, India,
        Summer 2013:
        Derived rules from the CHAID decision tree analysis to
        efficiently create appropriate customer segments for selling
        specific electronic products and rules to enhance cross-sale
        opportunities.
        \item eClerx, Mumbai, India, Summer 2012: Developed training decks on linear regression and cluster analysis for onboarding new employees. 
  
    \end{itemize}

  \resumeSubHeadingListEnd

\section*{Personal}

\begin{itemize}
\item Born in India to parents working as teachers and school administrators. 
\item I can speak 4 languages including 
English, Telugu, Hindi and Bengali and can write in three. 
\item Love Traveling and Sports. 
\item Have been to 36/63 National Parks in the U.S and traveled across the U.S. thrice. 
\item Had club membership with Badminton, Table-Tennis, Racquetball and Volleyball.
\item Always keen to learn more from any of mathematical, engineering or creative side. 

\end{itemize} 





%-------------------------------------------
\end{document}